\section*{Chapter \ref{ch:rc:rates-risk} --- Rates Risk}

\begin{solution}{exo:rc:rates-risk:1}
Parallel DV01 $50-120+80 = +10$ thousand dollars per basis point. P\&L
$50\times(-5) - 120\times2 + 80\times6 = -250-240+480 = -10$ thousand dollars.
\end{solution}

\begin{solution}{exo:rc:rates-risk:2}
They describe two different scenarios. The par ladder sums to the change for a
one-basis-point rise of every \emph{quote} with recalibration, which the
balancing swaps set to zero; the zero ladder sums to the change for a
one-basis-point rise of every \emph{\omterm{def:rc:curve-construction:pillar}{pillar} zero rate}, which moves each par
quote by slightly more than a basis point (the rates differ in compounding and
day count) and by different amounts, so the book's offsetting buckets no
longer cancel.
\end{solution}

\begin{solution}{exo:rc:rates-risk:3}
$100 - 98.2 = 1.8\%$ of the variance. The \omterm{def:rc:rates-risk:factors}{level factor}'s daily standard
deviation is $\sqrt{\lambda_1} = 13.6$ basis points of score, a move of about
$13.6\times0.40\approx5$ basis points at the belly of the curve.
\end{solution}

\begin{solution}{exo:rc:rates-risk:4}
$367+713+1\,531+3\,205+5\,120+74\,016 = 84\,952$, against a par ladder of
$83\,778$: the zero ladder is larger by the factor $1.014$ by which a
one-basis-point shift of the continuously compounded ACT/365 zero rates moves
the annual ACT/360 par rate. The shorter \omterm{def:rc:curve-construction:pillar}{pillars} appear because raising an early
zero rate lowers the value of the fixed coupons the payer owes; in quote space
recalibration absorbs that effect into the ten-year \omterm{def:rc:curve-construction:pillar}{pillar}.
\end{solution}

\begin{solution}{exo:rc:rates-risk:5}
Key-rate shifts are defined on the zero curve itself, whatever instruments and
interpolation built it; a par ladder depends on the instrument set and, off
\omterm{def:rc:curve-construction:pillar}{pillars}, on the interpolation (chapter~1). Send key rates (or a regulator's own
buckets) for comparability across desks; give the trader the par ladder, whose
buckets are directly the notionals of hedging par swaps.
\end{solution}

\begin{solution}{exo:rc:rates-risk:6}
The hedge minimises variance, not buckets: the seven- and ten-year risks have
opposite signs on highly correlated maturities, so together they carry little
variance (the ten-year roughly offsets the seven-year); hedging them would cost
notional and bid--offer for almost no reduction in risk. They are a small
seven-against-ten-year spread position.
\end{solution}

\begin{solution}{exo:rc:rates-risk:7}
Best single swap: the twenty-year, which leaves 86.5\% of the variance. Best
four: two, seven, twenty and thirty years, which leave 2.85\%, against 4.2\%
for the best three: the fourth swap buys little.
\end{solution}

\begin{solution}{exo:rc:rates-risk:8}
Being neutral to the first three components leaves 68.9\% of this book's
variance: the three factors describe how the whole curve moves, not the
twenty-against-thirty-year spread where the book's risk sits, which the higher
components carry. Factor neutrality removes factor risk only; the residual must
be measured on the full covariance, or hedged in the buckets where it lies.
\end{solution}

\begin{solution}{pb:rc:rates-risk:1}
\textbf{1.} $g^\top\delta q = -\text{USD}~3\,000\,000$.
\textbf{2.} $-\text{USD}~3\,077\,970$; the $-77\,970$ difference is second order
(\cref{ex:rc:rates-risk:gamma}).
\textbf{3.} Thirty years, $+958\,308\times9 = +8.62$ million, and twenty years,
$-1\,135\,816\times7 = -7.95$ million; then two years ($-2.76$ million) and seven
($-2.00$).
\textbf{4.} A steepening (short yields down, long up) with a twist near ten years;
the book lost on the two-year and seven-year buckets, which are long rates
there, and its long-end gains and losses nearly cancelled.
\textbf{5.} It measures only the sum of the buckets, the response to a move no one
saw that day.
\textbf{6.} $+21.2$.
\textbf{7.} $-\text{USD}~3.09$ million.
\textbf{8.} $21.2/4.87 = 4.3$ standard deviations: a rare slope day.
\textbf{9.} Moves of individual maturities against their neighbours: the book's
large, opposite twenty- and thirty-year buckets make them matter.
\textbf{10.} Roughly for the slope part; not for the rest, since the book's risk
outside three factors is large (exercise~8).
\textbf{11.} USD~1\,288\,843 per day.
\textbf{12.} 58.3\%.
\textbf{13.} Receive fixed on USD~1.196 billion of two-year, pay fixed on USD~801
million of twenty-year, receive fixed on USD~515 million of thirty-year.
\textbf{14.} USD~264\,554 per day ($\sqrt{0.042}\times1\,288\,843$).
\textbf{15.} $-\text{USD}~315\,044$ (first order): the hedge removed nine tenths of
the loss.
\textbf{16.} Each hedge costs bid--offer and balance sheet, and zeroing every bucket
removes offsets that historical moves make almost riskless; a variance
criterion hedges what actually moves.
\textbf{17.} Monthly or weekly, on a window long enough for stable estimates
(several years) but weighted towards recent data, and checked against the
regime (policy on hold, cutting or hiking); stress the hedge on other windows.
\textbf{18.} The ladder, the factor exposures, the residual variance after the
standing hedges, the P\&L on a set of historical twist days, and the gamma.
\textbf{19.} Named result: \emph{the twisted day} costs the DV01-flat book USD~3.00
million at first order (3.08 by full revaluation); receiving two-year, paying
twenty-year and receiving thirty-year swaps leaves 4.2\% of its daily variance.
\textbf{20.} It is one direction out of many, and the curve almost never moves in
exactly that direction.
\end{solution}

\begin{solution}{iq:rc:rates-risk:1}
Only against parallel moves. Look at the bucketed ladder: a zero sum can hide
large offsetting buckets, a steepener or a butterfly, which lose on the moves
that actually happen. Ask for the ladder, the factor exposures and a set of
historical twist scenarios.
\iqlookfor{the difference between a scalar and a vector risk.}
\end{solution}

\begin{solution}{iq:rc:rates-risk:2}
With the Jacobian $J = \partial(\text{model quotes})/\partial(\text{pillar
zeros})$ from the calibration: $g_z = J^\top g_q$, and $g_q =
(J^\top)^{-1}g_z$. Or by brute force: bump quotes and recalibrate, or bump
\omterm{def:rc:curve-construction:pillar}{pillars} without recalibration. Check that both routes agree.
\iqlookfor{the chain rule through the calibration.}
\end{solution}

\begin{solution}{iq:rc:rates-risk:3}
Level (all maturities the same sign), slope (short against long) and curvature
(belly against wings); on daily US Treasury changes 2016--2026 they explain
about 85\%, 11\% and 2\%, 98\% together. The shares depend on the market and
the period (the front end behaves differently near policy moves).
\iqlookfor{the shapes and orders of magnitude, and that they are estimates.}
\end{solution}

\begin{solution}{iq:rc:rates-risk:4}
Compute the book's ladder and the covariance of quote moves; for each
candidate triple of liquid swaps solve the minimum-variance hedge and keep the
triple with the least residual variance, with a check that it holds on another
window. Prefer liquid \omterm{def:rc:curve-construction:pillar}{pillars}, penalise large notionals. Here the textbook
two, ten and thirty years leave 58\% and two, twenty and thirty leave 4\%.
\iqlookfor{a criterion, not a rule of thumb.}
\end{solution}

\begin{solution}{iq:rc:rates-risk:5}
Second-order terms: the gamma matrix, including \omterm{def:rc:rates-risk:crossgamma}{cross-gammas} between buckets,
and for option books vanna and volga. Compute it by central differences of the
ladder or of the value, and add $\tfrac12\delta q^\top\Gamma\delta q$; if a
residual remains, look for missing risk factors (basis, volatility) or stale
inputs.
\iqlookfor{Taylor expansion to second order, then missing factors.}
\end{solution}

\begin{solution}{iq:rc:rates-risk:6}
When the book's risk lies outside the first factors (spread positions between
neighbouring maturities); when the estimation window's regime differs from
today's (zero rates, then hikes); when the factors are unstable in sign or
order; and on stress days, when correlations change. Always report the residual
on the full covariance.
\iqlookfor{PCA as a model with its own risk.}
\end{solution}
